%======================================================================
\section{Introduction}
\label{sec:intro}
%======================================================================

The adoption of Large Language Models (LLMs) for code generation has surged, with tools such as ChatGPT, GitHub Copilot, and Claude becoming integral to software development workflows~\cite{copilot_eval, fan2023large, hou2024large}. A 2023 GitHub survey reported that 92\% of developers use AI coding tools, with 70\% citing productivity gains~\cite{github_copilot}. This trend extends to the domain of smart contract development, where Solidity contracts govern billions of dollars in decentralized finance (DeFi) protocols~\cite{werner2022sok, zhou2023sok}. Unlike conventional software, vulnerabilities in smart contracts can lead to immediate, irreversible financial losses---as evidenced by over \$3.8 billion lost to smart contract exploits in 2022 alone~\cite{chainalysis2023}.

The convergence of LLM code generation and smart contract development creates a particularly dangerous intersection. Smart contracts are immutable once deployed: unlike web applications or traditional software, post-deployment patching is impossible without costly proxy-based migration~\cite{chen2020defining}. Furthermore, smart contracts operate in an adversarial environment where any vulnerability is immediately exploitable by any observer of the public blockchain. This raises a critical question: \emph{when developers use LLMs to generate smart contracts, what security vulnerabilities are systematically introduced, and how do these compare across models and prompting strategies?}

Despite the high-stakes nature of smart contracts, the security implications of LLM-generated Solidity code remain largely unexplored. Prior work has examined LLM-assisted vulnerability \emph{detection}~\cite{gptscan, smartllmsentry, david2023you}, but no study has systematically investigated vulnerabilities \emph{introduced} by LLMs during contract generation. This gap is particularly concerning given three converging factors: (1)~developers increasingly rely on LLM output with minimal review---Perry et al.~\cite{perry2023users} found that developers using AI assistants produce less secure code while believing it to be more secure; (2)~LLMs may encode vulnerable patterns from training data, including the vast corpus of unaudited contracts on Etherscan~\cite{schuster2021you}; and (3)~adversarial prompts can systematically induce security flaws in LLM outputs~\cite{kang2024exploiting}.

The problem is further compounded by the rapid expansion of cross-chain bridge protocols, which have become the highest-value attack surface in blockchain security, with over \$2.5 billion lost to bridge exploits~\cite{bridge_hacks}. If LLMs systematically generate vulnerable bridge contracts, the financial impact could be catastrophic---a single bridge exploit (Ronin, \$625M) exceeds the total losses from thousands of individual contract vulnerabilities.

\subsection{Research Questions}

We formulate six research questions to systematically investigate the security landscape of LLM-generated smart contracts:

\begin{itemize}[leftmargin=*]
    \item \textbf{RQ1 (Prevalence):} What is the overall prevalence of security vulnerabilities in LLM-generated smart contracts, and how does it compare to the vulnerability rates reported for human-written contracts?
    \item \textbf{RQ2 (Cross-LLM):} Do statistically significant differences exist in the security posture of contracts generated by different LLMs, and which models produce the most and least secure code?
    \item \textbf{RQ3 (Taxonomy):} What types of vulnerabilities are most prevalent in LLM-generated smart contracts, and can we construct a comprehensive CWE-mapped taxonomy?
    \item \textbf{RQ4 (Adversarial):} How do adversarial prompts---those that exploit LLM biases toward gas optimization, simplicity, or speed---affect vulnerability rates compared to standard prompts?
    \item \textbf{RQ5 (Patterns):} Do LLMs exhibit distinct vulnerability patterns that are absent in human-written code, and can these patterns be characterized for targeted detection?
    \item \textbf{RQ6 (Human Comparison):} How do LLM-generated contracts compare to audited, human-written contracts across all vulnerability dimensions?
\end{itemize}

\subsection{Contributions}

In this paper, we present the first large-scale empirical study of security vulnerabilities in LLM-generated smart contracts. Our study makes the following contributions:

\begin{itemize}[leftmargin=*]
    \item \textbf{Large-Scale Dataset.} We generate 1,800 smart contracts across six LLMs (GPT-4o, Claude~3.5 Sonnet, Gemini~2.5 Flash-Lite, DeepSeek Coder~V2, Qwen~2.5 Coder~32B, CodeLlama-34B), covering 20 contract categories at three complexity levels. This is the largest dataset of LLM-generated smart contracts with security labels to date.
    \item \textbf{Multi-Tool Analysis Pipeline.} We employ three complementary security analysis tools---Slither, Mythril, and Semgrep---with cross-tool deduplication and manual validation, achieving approximately 92\% vulnerability detection coverage.
    \item \textbf{Vulnerability Taxonomy.} We construct a CWE-mapped taxonomy of vulnerability types specific to LLM-generated smart contracts, identifying three novel patterns (phantom guards, inherited vulnerabilities, context window degradation).
    \item \textbf{Adversarial Prompt Framework.} We design eight categories of adversarial prompts that exploit documented LLM biases (e.g., gas optimization, simplicity). Contrary to expectations, adversarial prompts do \emph{not} increase vulnerability rates---we show vulnerabilities are intrinsic to LLM code generation.
    \item \textbf{Cross-LLM Comparison.} We provide statistically rigorous comparisons across LLMs using non-parametric tests with effect size reporting, revealing significant differences in security posture ($p < 0.001$).
\end{itemize}

Our results have direct implications for developers, LLM providers, and the security community. We release our full dataset, analysis scripts, and tooling to support reproducibility and future research.

\textbf{Paper Organization.} Section~\ref{sec:background} provides background on smart contract security and LLM code generation. Section~\ref{sec:threat} defines our threat model. Section~\ref{sec:methodology} describes our methodology. Sections~\ref{sec:results} and~\ref{sec:category} present our findings. Section~\ref{sec:discussion} discusses implications and limitations. Section~\ref{sec:related} surveys related work. Section~\ref{sec:conclusion} concludes.
